\documentclass[12pt]{article}
\usepackage[utf8]{inputenc}
\usepackage[T1]{fontenc}
\usepackage{hyperref}
\title{Tera Language Resources for K-12 Curriculum Development}
\author{Compiled by Claude Sonnet 5 (claude-sonnet-5) for Andr\'as Kornai}
\date{}
\begin{document}
\maketitle
\section{Introduction}
Tera (ISO 639-3 ttr; Glottocode tera1251) is a Chadic (Afro-Asiatic) language of Gombe and Borno States, Nigeria. Speaker estimates range from about 100,000 to about 200,000. The community also uses the name Nyimatli. Only a few resources turned up in search, and no school curriculum materials were found.
\section{Primary Educational Resources}
No textbooks or school primers were found in search. The orthography literature mentions a reading and writing guide, \emph{Reading \& Writing Nyimatli} (definitive edition 2015), from the Nyimatli Language Project. I could not confirm the details of this guide from the search results.
\begin{itemize}
\item Chapter on Tera in a scholarly volume (orthography background): \url{https://api.equinoxpub.com/chapter/show/44120}
\end{itemize}
\section{Linguistic and Reference Materials}
\begin{itemize}
\item Tench, ``Tera'' (Illustrations of the IPA, Journal of the International Phonetic Association): \url{https://resolve.cambridge.org/core/services/aop-cambridge-core/content/view/EBA55DEC3AD50C945BC3DC58C78DE246/S0025100307002952a.pdf/tera.pdf}
\item LINGUIST List notice of the \emph{Tera-English-Hausa Dictionary} (Maimuna Magaji, 2018): \url{https://linguistlist.org/issues/26/648/}
\item LINGUIST List thread on Tera transcription: \url{https://listserv.linguistlist.org/pipermail/ilat/2005-July/001649.html}
\item Wikipedia overview: \url{https://en.wikipedia.org/wiki/Tera_language}
\item Rosetta Panglossia entry: \url{https://rosettapanglossia.longnow.org/index.php?title=Ttr}
\end{itemize}
\section{Religious and Community Resources}
The Gospel of John (1930) was the first printed text. The New Testament, \emph{Mewar Alqawarang}, was published in 2016 by the Nigerian Bible Translation Trust. These facts come from search summaries. I did not see a direct link to the scripture text.
\begin{itemize}
\item Joshua Project profile: \url{https://joshuaproject.net/people_groups/15440}
\end{itemize}
\section{Digital and Audio Resources}
No dedicated audio, video or software resources were found in search.
\section{Curriculum Development Context}
Tera is not known to be taught in schools. The orthography is recent (workshops in 2004 and 2015) and does not mark tone. The dictionary is trilingual (Tera, English, Hausa), which could support bilingual materials. Primers and readers would have to be newly developed with the community.
\end{document}